\documentclass[12pt, a4paper]{article}

% ============================================================
%  PACKAGES
% ============================================================
\usepackage[margin=2.5cm]{geometry}
\usepackage{amsmath, amssymb, amsfonts}
\usepackage{graphicx}
\usepackage{hyperref}
\usepackage{xcolor}
\usepackage{titlesec}
\usepackage{fancyhdr}
\usepackage{enumitem}
\usepackage{booktabs}
\usepackage{array}
\usepackage{listings}
\usepackage{caption}
\usepackage{subcaption}
\usepackage{parskip}
\usepackage{setspace}
\usepackage{tcolorbox}

% ============================================================
%  COLOURS & STYLING
% ============================================================
\definecolor{primary}{RGB}{59,130,246}      % blue-500
\definecolor{accent}{RGB}{16,185,129}       % emerald-500
\definecolor{warning}{RGB}{245,158,11}      % amber-500
\definecolor{danger}{RGB}{251,113,133}      % rose-400
\definecolor{purple}{RGB}{168,85,247}       % purple-500
\definecolor{codebg}{RGB}{23,23,23}
\definecolor{neutral}{RGB}{115,115,115}

\hypersetup{
  colorlinks=true,
  linkcolor=primary,
  urlcolor=primary,
  citecolor=accent
}

\titleformat{\section}{\Large\bfseries\color{primary}}{}{0em}{\thesection\quad}[\titlerule]
\titleformat{\subsection}{\large\bfseries\color{accent}}{}{0em}{\thesubsection\quad}
\titleformat{\subsubsection}{\normalsize\bfseries\color{warning}}{}{0em}{\thesubsubsection\quad}

\pagestyle{fancy}
\fancyhf{}
\rhead{\textcolor{neutral}{\small Stochastic Measurement \& Sensor Fusion}}
\lhead{\textcolor{neutral}{\small BECE207L --- Random Processes}}
\cfoot{\thepage}
\renewcommand{\headrulewidth}{0.4pt}

\lstset{
  basicstyle=\ttfamily\small\color{white},
  backgroundcolor=\color{codebg},
  breaklines=true,
  frame=single,
  rulecolor=\color{neutral},
  keywordstyle=\color{primary},
  commentstyle=\color{neutral},
  stringstyle=\color{accent},
  numbers=left,
  numberstyle=\tiny\color{neutral},
  numbersep=8pt,
  tabsize=2,
  captionpos=b
}

% ============================================================
%  DOCUMENT
% ============================================================
\begin{document}

% ====================================================
%  TITLE PAGE
% ====================================================
\begin{titlepage}
  \centering
  \vspace*{1cm}

  {\Large\bfseries Vellore Institute of Technology}\\[0.3cm]
  {\large School of Electronics Engineering}\\[0.3cm]
  {\large BECE207L --- Random Processes}\\[2cm]

  \hrule height 2pt
  \vspace{0.6cm}
  {\Huge\bfseries\color{primary} Stochastic Measurement \&\\[0.4cm] Sensor Fusion}\\[0.4cm]
  {\LARGE\color{accent} An Interactive Simulation Dashboard}
  \vspace{0.6cm}
  \hrule height 2pt

  \vspace{1.5cm}
  {\large\bfseries A Course Project Report}\\[0.3cm]
  {\normalsize Submitted in partial fulfilment of the requirements of BECE207L}

  \vfill

  \begin{tabular}{ll}
    \toprule
    \textbf{Student} & \textbf{Registration Number} \\
    \midrule
    S Malaiesun     & 24BEC1578 \\
    T S Vishal Babu & 24BEC1575 \\
    J Saran         & 24BEC1081 \\
    \bottomrule
  \end{tabular}

  \vspace{1cm}
  {\normalsize\textbf{Faculty Guide:} Prof. Kalaivanan K}\\[0.3cm]
  {\normalsize\textbf{Date:} April 2026}
\end{titlepage}

% ====================================================
%  TABLE OF CONTENTS
% ====================================================
\tableofcontents
\newpage
\onehalfspacing

% ====================================================
%  ABSTRACT
% ====================================================
\begin{tcolorbox}[colback=primary!5!white, colframe=primary, title=\textbf{Abstract}]
  This report documents the design, mathematical theory, and implementation of an
  interactive web-based simulation platform built for the BECE207L Random Processes
  course. The platform provides a unified, seven-module environment for visualising and
  experimenting with the full signal-processing pipeline encountered in modern sensor
  systems: from the foundational mathematics of probability distributions, through
  additive white Gaussian noise (AWGN) and Brownian random-walk drift, to diagnostic
  signal analysis, optimal state estimation via the discrete Kalman filter, adaptive
  noise cancellation via the Least Mean Squares (LMS) algorithm, and finally
  variance-weighted sensor fusion. Every algorithm is implemented from scratch in
  TypeScript without any external numerics library, and every result is rendered in
  real-time interactive charts. The platform targets students and practising engineers
  who need an intuitive yet mathematically rigorous laboratory for stochastic signal
  processing.
\end{tcolorbox}

\vspace{1em}

% ====================================================
\section{Introduction}
% ====================================================

\subsection{Motivation}

All physical measurements are corrupted by noise. Whether it is a MEMS inertial
measurement unit (IMU) on an aircraft, a temperature sensor on an embedded controller,
or a radar return from a target, the signal that reaches a digital processor is never a
clean replica of the physical quantity being observed.  Understanding the statistical
nature of that corruption is not merely academic: it is the prerequisite to designing
filters and estimators that can recover the true quantity.

The Random Processes course (BECE207L) covers the theoretical foundations of stochastic
signals --- probability distributions, wide-sense stationarity, power spectral density,
autocorrelation, the Kalman filter, and adaptive algorithms. However, the gap between
textbook equations and observable behaviour is often large. This project bridges that
gap by providing a live, parameter-driven simulation laboratory.

\subsection{Project Objectives}

\begin{enumerate}[leftmargin=*, label=\arabic*.]
  \item Implement, from first principles in TypeScript, every major algorithm taught in
        the course.
  \item Provide real-time interactive visualisations so that the effect of changing a
        single parameter (e.g.\ noise standard deviation, Kalman process noise $Q$,
        LMS learning rate $\mu$) is immediately visible.
  \item Structure the modules in a logical pedagogical pipeline, starting from the
        mathematics of randomness and ending at hardware-level sensor fusion.
  \item Allow the user to upload real sensor CSV data for Kalman filtering and signal
        analysis, verifying the algorithms on non-synthetic data.
\end{enumerate}

\subsection{Technology Stack}

The application is a Next.js 15 web application written in TypeScript. The UI is built
with Tailwind CSS and shadcn/ui components. All interactive charts use Recharts.
Animations are driven by Framer Motion. The entire numerics layer resides in six
TypeScript modules inside the \texttt{/lib} directory --- no external maths library is
used.

\begin{table}[h!]
  \centering
  \caption{Technology Stack Overview}
  \begin{tabular}{lll}
    \toprule
    \textbf{Layer} & \textbf{Technology} & \textbf{Purpose} \\
    \midrule
    Framework     & Next.js 15 (App Router) & SSR / client routing \\
    Language      & TypeScript              & Type-safe numerics and UI \\
    Styling       & Tailwind CSS            & Responsive dark-mode layout \\
    Components    & shadcn/ui               & Sliders, Cards, Tabs, Buttons \\
    Charting      & Recharts                & Real-time line, bar, area charts \\
    Animation     & Framer Motion           & Cinematic transitions \\
    Numerics      & Custom \texttt{/lib}    & All stochastic algorithms \\
    \bottomrule
  \end{tabular}
\end{table}

% ====================================================
\section{System Architecture}
% ====================================================

\subsection{Module Structure}

The codebase is divided into two layers:

\begin{itemize}
  \item \textbf{Engine layer} (\texttt{/lib}): Pure TypeScript functions that perform
        numerical computation. They have no UI dependencies and can be unit-tested
        independently.
  \item \textbf{Simulation layer} (\texttt{/components/simulations}): React components
        that expose the engine functions to the user via sliders, CSV uploads, and charts.
\end{itemize}

\begin{table}[h!]
  \centering
  \caption{Engine Modules and Their Responsibilities}
  \begin{tabular}{>{\ttfamily}ll}
    \toprule
    \textbf{File} & \textbf{Responsibility} \\
    \midrule
    stochasticEngine.ts & Gaussian RNG (Box-Muller), time-series generation, RMSE/variance \\
    driftEngine.ts      & Brownian random-walk (integrated Gaussian steps) \\
    filterEngine.ts     & 1-D discrete Kalman filter, CSV parser \\
    lmsEngine.ts        & LMS adaptive FIR filter (stochastic gradient descent) \\
    fusionEngine.ts     & Variance-weighted sensor fusion (Kalman update logic) \\
    analysisEngine.ts   & SNR, PSD (DFT), autocorrelation, WSS test, moments, process ID \\
    \bottomrule
  \end{tabular}
\end{table}

\subsection{Application Flow}

The root \texttt{page.tsx} implements a two-state machine:

\begin{enumerate}[label=\arabic*.]
  \item \textbf{Landing View} --- A cinematic hero page that introduces the project,
        lists contributors, and summarises all seven modules.
  \item \textbf{Dashboard View} --- A seven-tab interface where each tab shows (a)~a
        theory panel explaining the mathematics, and (b)~the live simulation widget.
\end{enumerate}

Framer Motion's \texttt{AnimatePresence} drives blur/scale transitions between pages
and between active tabs, giving the interface a fluid, aerospace-grade aesthetic.

% ====================================================
\section{Mathematical Foundations \& Module Theory}
% ====================================================

The seven modules are ordered as a logical pipeline. Each module builds on the concepts
established by the previous one.

% --------------------------------------------------
\subsection{Module 1 --- Statistical Distributions}
% --------------------------------------------------

\subsubsection{The Core Problem}

Computers generate \emph{uniform} pseudo-random numbers: each value in $[0, 1)$ is
equally likely. Physical noise, however, is not uniform --- it follows a
\textbf{Gaussian (Normal) distribution}, where extreme values are exponentially rare.
Before any sensor can be modelled, we must be able to generate Gaussian samples.

\subsubsection{The Box-Muller Transform}

Given two independent uniform random variables $U_1, U_2 \sim \mathcal{U}(0,1)$, the
transform produces two independent standard normal samples:

\begin{equation}
  Z_0 = \sqrt{-2\ln U_1}\;\cos(2\pi U_2), \qquad
  Z_1 = \sqrt{-2\ln U_1}\;\sin(2\pi U_2)
  \label{eq:boxmuller}
\end{equation}

A general Gaussian sample with mean $\mu$ and standard deviation $\sigma$ is then:
\begin{equation}
  X = Z_0 \cdot \sigma + \mu
\end{equation}

The implementation in \texttt{stochasticEngine.ts}:
\begin{itemize}
  \item Guards against $\ln(0) \to -\infty$ by clamping $U_1 \geq 10^{-7}$.
  \item Uses only $Z_0$ (one sample per call) for simplicity.
\end{itemize}

\subsubsection{What the Simulation Demonstrates}

The user toggles between Uniform and Gaussian sampling and increases the sample count
$N$. By the \textbf{Law of Large Numbers}, the empirical histogram converges to:
\begin{itemize}
  \item A flat rectangular block for Uniform.
  \item A symmetric bell curve $\mathcal{N}(\mu, \sigma^2)$ for Gaussian.
\end{itemize}

% --------------------------------------------------
\subsection{Module 2 --- Additive White Gaussian Noise (AWGN)}
% --------------------------------------------------

\subsubsection{Definition}

AWGN is the canonical model for sensor noise. A sensor measurement $Z(t)$ of a true
signal $X$ is modelled as:

\begin{equation}
  Z(t) = X + \beta + \varepsilon(t), \qquad \varepsilon(t) \sim \mathcal{N}(0, \sigma^2)
  \label{eq:awgn}
\end{equation}

where $\beta$ is a \emph{constant bias} (systematic offset) and $\varepsilon(t)$ is
zero-mean Gaussian noise (random jitter). The three key statistical descriptors are:

\begin{align}
  \text{Mean Error:} \quad & \mu_e = \frac{1}{N}\sum_{t=1}^{N} e(t) \approx \beta \\
  \text{Variance:}   \quad & \sigma_e^2 = \frac{1}{N}\sum_{t=1}^{N}(e(t) - \mu_e)^2 \approx \sigma^2 \\
  \text{RMSE:}       \quad & \text{RMSE} = \sqrt{\frac{1}{N}\sum_{t=1}^{N} e(t)^2}
\end{align}

\subsubsection{Why AWGN is Called ``White''}

White light contains all visible frequencies at equal intensity. Analogously, AWGN has
a \emph{flat} Power Spectral Density (PSD) --- it contributes equal power at every
frequency. This property is verified in Module 4.

\subsubsection{What the Simulation Demonstrates}

The user adjusts the true signal $X$, bias $\beta$, noise $\sigma$, and sample count
$N$. The left panel displays live statistics; the right panel shows (a)~the time series
of the true vs.\ measured signal, and (b)~a histogram of the error samples. As
$\sigma \to 0$, the histogram collapses to a delta function at $\beta$. As
$N \to \infty$, the histogram shape converges to a bell curve --- confirming the
Central Limit Theorem.

% --------------------------------------------------
\subsection{Module 3 --- Brownian Motion \& Random Walk}
% --------------------------------------------------

\subsubsection{The IMU Drift Problem}

Inertial Measurement Units (IMUs) integrate acceleration to produce velocity, and
integrate velocity to produce position. Even a tiny, zero-mean noise sample
$\varepsilon(t)$ produces a \emph{cumulative} error when integrated: the error
performs a \textbf{Random Walk} (Brownian motion) and wanders infinitely away from the
true state over time. This is a fundamentally different failure mode from AWGN.

\subsubsection{Mathematical Model}

The discrete random-walk bias evolves as:

\begin{equation}
  b(t) = b(t-1) + \delta(t), \qquad \delta(t) \sim \mathcal{N}(0, \sigma_d^2)
  \label{eq:randomwalk}
\end{equation}

The total sensor measurement at time $t$ is:

\begin{equation}
  Z(t) = X + b(t) + \varepsilon(t)
\end{equation}

where $b(t)$ is the \emph{accumulated} drift and $\varepsilon(t)$ is the high-frequency
jitter. Because each step $\delta(t)$ is added cumulatively, the variance of $b(t)$
grows linearly with time:

\begin{equation}
  \text{Var}[b(t)] = t \cdot \sigma_d^2
\end{equation}

This unbounded growth makes the random walk a \emph{non-stationary} process.

\subsubsection{What the Simulation Demonstrates}

When the jitter is zeroed and drift severity is increased, the rose-coloured sensor
line wanders further from the green true-signal line with every time step. This
visually proves that IMUs used in standalone mode (without GPS or magnetometer
correction) will eventually lose spatial coherence --- motivating the sensor fusion
approach in Module 7.

% --------------------------------------------------
\subsection{Module 4 --- DSP Signal Diagnostics}
% --------------------------------------------------

\subsubsection{Purpose}

Before applying a filter, an engineer must \emph{diagnose} the noise: What statistical
process generated it? Is it stationary? What is the pre-filter Signal-to-Noise Ratio?
This module automates that diagnosis.

\subsubsection{Signal-to-Noise Ratio (SNR)}

\begin{equation}
  \text{SNR} = 10\log_{10}\!\left(\frac{P_\text{signal}}{P_\text{noise}}\right) \quad \text{[dB]}
  \label{eq:snr}
\end{equation}

where $P = \frac{1}{N}\sum x_i^2$ is the mean power. A higher SNR means more of the
measured energy is true signal; a lower SNR means the noise is dominant.

\subsubsection{Autocorrelation Function}

The autocorrelation at lag $k$ measures how similar a signal is to a time-shifted
version of itself:

\begin{equation}
  R_{xx}(k) = \frac{1}{N}\sum_{i=0}^{N-k-1} x(i)\, x(i+k)
  \label{eq:autocorr}
\end{equation}

For white noise, $R_{xx}(k) = 0$ for all $k \neq 0$ (no memory). For coloured
(Markov) noise, $R_{xx}(k)$ decays exponentially. For a random walk (Wiener process),
$R_{xx}(k)$ decays extremely slowly, retaining high correlation at large lags.

\subsubsection{Power Spectral Density (PSD)}

The PSD is computed via the Discrete Fourier Transform (DFT):

\begin{equation}
  S(k) = \frac{1}{N}\left|\sum_{n=0}^{N-1} x(n)\, e^{-j2\pi kn/N}\right|^2
  \label{eq:psd}
\end{equation}

White noise: flat PSD. Coloured noise: sloping PSD. The implementation uses
a direct $O(N^2)$ DFT (sufficient for pedagogical $N \leq 200$).

\subsubsection{Wide-Sense Stationarity (WSS) Test}

A process is Wide-Sense Stationary if its mean and variance are constant over time.
The implementation splits the data in half and compares the two halves:

\begin{align}
  \Delta\mu &= |\mu_1 - \mu_2| < \sqrt{\max(\sigma_1^2, \sigma_2^2)} \times 0.8 \\
  r_\sigma  &= \frac{\max(\sigma_1^2, \sigma_2^2)}{\min(\sigma_1^2, \sigma_2^2)} < 3.0
\end{align}

Both conditions must hold for the WSS flag to pass. The thresholds are
\emph{scale-invariant} (derived from the data's own variance), avoiding the hardcoded
numerical threshold that would fail on differently-scaled inputs.

\subsubsection{Statistical Moments \& Process Identification}

Higher statistical moments characterise the shape of the probability distribution:

\begin{align}
  \text{Skewness:} \quad & \gamma_1 = \frac{1}{N\sigma^3}\sum(x_i - \mu)^3 \\
  \text{Kurtosis:} \quad & \gamma_2 = \frac{1}{N\sigma^4}\sum(x_i - \mu)^4
\end{align}

A Gaussian distribution has $\gamma_1 \approx 0$ and $\gamma_2 \approx 3$.
Deviations identify different random processes. The \texttt{identifyProcess()} function
implements the following decision tree:

\begin{table}[h!]
  \centering
  \caption{Automated Process Classification}
  \begin{tabular}{llll}
    \toprule
    \textbf{Condition} & \textbf{Classification} \\
    \midrule
    $\gamma_2 > 4.5$ & Poisson (Impulsive) \\
    $R(1)/R(0) > 0.9$ \emph{and} $R(10)/R(0) > 0.6$ & Wiener (Brownian Drift) \\
    $R(1)/R(0) > 0.3$ & Markov (Coloured Noise) \\
    Otherwise & Gaussian (White Noise) \\
    \bottomrule
  \end{tabular}
\end{table}

\subsubsection{Low-Pass Filter for SNR Improvement}

A causal moving-average low-pass filter (LPF) of window size $W$ is applied:

\begin{equation}
  y(t) = \frac{1}{\min(t+1,\, W)}\sum_{j=0}^{\min(t, W-1)} x(t-j)
\end{equation}

The SNR improvement (in dB) after filtering is reported as a net gain metric.

% --------------------------------------------------
\subsection{Module 5 --- 1-D Discrete Kalman Filter}
% --------------------------------------------------

\subsubsection{Motivation}

The Kalman filter is the \emph{optimal} linear estimator for Gaussian noise in the
least-squares sense. It outperforms a fixed moving average because it dynamically
allocates trust between its own prediction model and the incoming measurement,
weighted by their respective uncertainties.

\subsubsection{State-Space Model}

For a static (or slowly varying) 1-D scalar state $x_t$:

\begin{align}
  \text{Process model:}     \quad & x_t = x_{t-1} + w_t, \quad w_t \sim \mathcal{N}(0, Q) \\
  \text{Measurement model:} \quad & z_t = x_t + v_t,     \quad v_t \sim \mathcal{N}(0, R)
\end{align}

$Q$ is the \emph{process noise variance} (how much the true state can change between
steps) and $R$ is the \emph{measurement noise variance} (how noisy the sensor is).

\subsubsection{Algorithm}

The filter alternates between two steps at each time $t$:

\textbf{Prediction:}
\begin{align}
  \hat{x}_{t|t-1} &= \hat{x}_{t-1|t-1} \\
  P_{t|t-1}       &= P_{t-1|t-1} + Q
\end{align}

\textbf{Update:}
\begin{align}
  K_t             &= \frac{P_{t|t-1}}{P_{t|t-1} + R}              && \text{(Kalman Gain)} \\
  \hat{x}_{t|t}   &= \hat{x}_{t|t-1} + K_t(z_t - \hat{x}_{t|t-1}) && \text{(Updated Estimate)} \\
  P_{t|t}         &= (1 - K_t)\, P_{t|t-1}                         && \text{(Updated Covariance)}
\end{align}

\textbf{Interpretation of the Kalman Gain $K_t$:}
\begin{itemize}
  \item $K_t \to 1$: Trust the raw measurement almost entirely (large $R$ relative
        to $P$ means prior is uncertain).
  \item $K_t \to 0$: Trust the prediction almost entirely (small $R$ means sensor
        is very accurate).
\end{itemize}

\subsubsection{What the Simulation Demonstrates}

The user uploads a CSV file of raw sensor readings. The $R$ slider (``Sensor
Distrust'') is tuned live. Increasing $R$ produces heavy smoothing; decreasing $R$
allows the filter to track the raw data more aggressively. The visualisation overlays
the raw signal and the filtered estimate, allowing the user to tune $R$ until they
achieve the desired balance between noise rejection and signal tracking.

% --------------------------------------------------
\subsection{Module 6 --- LMS Adaptive Filter}
% --------------------------------------------------

\subsubsection{Motivation: The Failure of Static Filters}

The Kalman filter requires $Q$ and $R$ to be set before deployment. If the noise
profile of the environment changes (e.g.\ a new interference source appears), a
static filter degrades without any self-correction mechanism. An \emph{adaptive}
filter learns the statistics of the noise online and continuously updates its
own coefficients.

\subsubsection{FIR Filter Structure}

A Finite Impulse Response (FIR) filter of order $M$ computes its output as:

\begin{equation}
  \hat{y}(t) = \sum_{i=0}^{M-1} w_i \cdot x(t-i)
  \label{eq:fir}
\end{equation}

where $\mathbf{w} = [w_0, w_1, \ldots, w_{M-1}]$ is the weight vector and
$[x(t), x(t-1), \ldots, x(t-M+1)]$ is the delay buffer of recent noisy inputs.

\subsubsection{Least Mean Squares (LMS) Algorithm}

The LMS algorithm minimises the mean squared error between the filter output
$\hat{y}(t)$ and the desired signal $d(t)$ (the true sine wave) using
\emph{stochastic gradient descent}:

\begin{align}
  e(t)     &= d(t) - \hat{y}(t)                           && \text{(Instantaneous Error)} \\
  w_i(t+1) &= w_i(t) + 2\mu \cdot e(t) \cdot x(t-i)      && \text{(Weight Update)}
  \label{eq:lms}
\end{align}

where $\mu > 0$ is the \textbf{learning rate}. The gradient-descent step size is
$2\mu e(t) x(t-i)$, derived from $-\nabla_{\mathbf{w}} \mathbb{E}[e^2(t)]$.

\subsubsection{Stability Condition}

The LMS algorithm converges if and only if:

\begin{equation}
  0 < \mu < \frac{1}{M \cdot P_x}
\end{equation}

where $P_x$ is the input signal power and $M$ is the filter order. Setting $\mu$ too
large causes the algorithm to diverge catastrophically (visible in the simulation as
the filtered output exploding to $\pm\infty$).

\subsubsection{What the Simulation Demonstrates}

At $t = 0$ all weights are zero, so the filter output is a flat line. Over time the
cyan filtered line ``learns'' to extract the underlying $5\sin(0.1t)$ sine wave from
the red noisy input. The user can push $\mu$ above the stability limit to observe the
mathematical chaos --- a visceral demonstration of why the stability bound matters.

% --------------------------------------------------
\subsection{Module 7 --- Variance-Weighted Sensor Fusion}
% --------------------------------------------------

\subsubsection{The Hardware Solution to Noise}

When a single sensor cannot be trusted, the solution is to use multiple sensors. If
two sensors $A$ and $B$ independently observe the same true quantity $X$ with
Gaussian noise:

\begin{align}
  z_A = X + \varepsilon_A, \quad \varepsilon_A &\sim \mathcal{N}(0, \sigma_A^2) \\
  z_B = X + \varepsilon_B, \quad \varepsilon_B &\sim \mathcal{N}(0, \sigma_B^2)
\end{align}

\subsubsection{Optimal Fusion Formula}

The minimum-variance (statistically optimal) fused estimate is the
\emph{variance-weighted average}:

\begin{equation}
  \hat{X} = \frac{z_A \cdot \sigma_B^2 + z_B \cdot \sigma_A^2}{\sigma_A^2 + \sigma_B^2}
  \label{eq:fusion}
\end{equation}

This is exactly the \textbf{Kalman update step} in scalar form. The fused variance is:

\begin{equation}
  \sigma_\text{fused}^2 = \frac{\sigma_A^2 \cdot \sigma_B^2}{\sigma_A^2 + \sigma_B^2}
\end{equation}

which is always smaller than both $\sigma_A^2$ and $\sigma_B^2$ --- the fusion is
mathematically guaranteed to be at least as good as the best individual sensor.

\subsubsection{Geometric Interpretation}

A noisy sensor with high variance $\sigma^2$ receives low weight in the fusion
(the denominator term $\sigma_A^2$ in the numerator is small when $\sigma_A^2$ is
small). The result is that the fused estimate automatically ``hugs'' the more reliable
sensor, dynamically adjusted whenever the user changes either $\sigma$.

\subsubsection{What the Simulation Demonstrates}

Three lines are plotted: Sensor A (amber), Sensor B (blue), and the fused estimate
(purple). As the user increases $\sigma_A$ (making sensor A noisier), the purple
fused line migrates toward sensor B. When both are equal, the purple line bisects the
two. The simulation also accepts CSV uploads for two real sensor streams, fusing them
with user-specified noise parameters.

% ====================================================
\section{Implementation Details}
% ====================================================

\subsection{Core Random Number Generation}

\texttt{stochasticEngine.ts} implements the Box-Muller transform and exposes
\texttt{generateGaussian($\mu$, $\sigma$)} as the shared building block for all other
engines. Every noise source in the platform flows through this single function,
guaranteeing numerical consistency across modules.

\subsection{CSV Data Pipeline}

All modules that accept real-world data use a shared regex-based CSV parser that
extracts the first numeric token from each non-empty line:
\begin{verbatim}
  const match = line.match(/-?\d+(\.\d+)?/);
\end{verbatim}
This approach handles CSVs with arbitrary headers, multi-column data (ignoring extra
columns), and both LF and CRLF line endings.

\subsection{Reactive Simulation Architecture}

Each simulation component uses a \texttt{useEffect} hook whose dependency array
includes all parameter state variables. Any slider adjustment immediately triggers
a fresh simulation run and re-renders all charts without user intervention, providing
a ``live'' laboratory feel.

\subsection{WSS Threshold Upgrade}

A known engineering pitfall of stationary tests is the use of absolute, hardcoded
thresholds that fail on differently-scaled data. The platform uses a
scale-invariant mean-difference threshold derived from the signal's own standard
deviation:
\begin{equation}
  \theta_\mu = \sqrt{\max(\sigma_1^2, \sigma_2^2)} \times 0.8
\end{equation}
This passes correctly for unit-amplitude noise and equally for signals scaled by
$10^6$, unlike a fixed absolute threshold.

% ====================================================
\section{Results \& Observations}
% ====================================================

\subsection{AWGN Module}

Setting bias $\beta = 0$ and varying $\sigma$ demonstrates that:
\begin{itemize}
  \item The mean error $\mu_e \approx 0$ for all $\sigma$ (unbiased estimator).
  \item The variance $\sigma_e^2$ tracks $\sigma^2$ within Monte Carlo sampling error.
  \item The RMSE $\approx \sigma$ when bias is zero, confirming $\text{RMSE}^2 = \text{Var} + \text{Bias}^2$.
\end{itemize}

\subsection{Brownian Drift Module}

With white noise at zero:
\begin{itemize}
  \item The sensor error grows without bound, consistent with
        $\text{Var}[b(t)] = t \cdot \sigma_d^2$.
  \item Doubling $\sigma_d$ approximately doubles the visible spread at any fixed $t$.
  \item The WSS test in Module 4 correctly returns \textbf{FAIL} for drifting data,
        since the variance of the second half is much larger than the first.
\end{itemize}

\subsection{Kalman Filter Module}

On synthetic Gaussian data:
\begin{itemize}
  \item With $R$ small (trusting the sensor), the filtered output tracks the raw
        signal closely with minimal smoothing.
  \item With $R$ large, the filtered output is heavily smoothed, converging toward the
        prior model (a constant).
  \item The optimal $R$ is that which minimises the residual RMS against the known
        ground truth.
\end{itemize}

\subsection{LMS Adaptive Filter Module}

\begin{itemize}
  \item At low $\mu$ (e.g.\ 0.001), the filter converges slowly but stably.
  \item At the recommended $\mu = 0.01$, the cyan line visually extracts the sine
        wave within approximately 50 samples.
  \item Pushing $\mu > 0.05$ causes visible numerical instability; $\mu > 0.1$
        produces divergence to $\pm\infty$, confirming the theoretical stability bound.
\end{itemize}

\subsection{Sensor Fusion Module}

\begin{itemize}
  \item The fused variance is always lower than the worse sensor, as proven by
        Equation \eqref{eq:fusion}.
  \item When $\sigma_A = \sigma_B$, the fused line runs exactly through the average
        of the two sensors.
  \item The fused line ``clings'' to whichever sensor has the lower $\sigma$, which
        is the mathematically optimal behaviour.
\end{itemize}

% ====================================================
\section{Conclusion}
% ====================================================

This project successfully implements a complete stochastic signal-processing pipeline
from first principles, covering:

\begin{enumerate}[label=\arabic*.]
  \item The mathematical transformation of uniform machine randomness into
        physically-meaningful Gaussian distributions (Box-Muller).
  \item The modelling of AWGN as the fundamental corruption mechanism in all sensors.
  \item The modelling of Brownian random-walk drift as a cumulative, non-stationary
        failure mode in integrating sensors such as IMUs.
  \item A full DSP diagnostics suite computing SNR, PSD, autocorrelation, the WSS
        test, and higher statistical moments for automatic process classification.
  \item A 1-D discrete Kalman filter that demonstrates the predict-update cycle and
        the effect of the process-noise/measurement-noise ratio on estimation quality.
  \item An LMS adaptive FIR filter that demonstrates online stochastic gradient descent
        and the critical importance of the learning-rate stability condition.
  \item Variance-weighted sensor fusion implementing the core Kalman update logic for
        multi-sensor environments.
\end{enumerate}

The interactive web platform makes the abstract mathematics of random processes
immediately observable and experimentally verifiable, serving as both a pedagogical
tool and a practical engineering sandbox.

\vspace{2em}

\begin{tcolorbox}[colback=accent!5!white, colframe=accent, title=\textbf{Key Insight}]
  Every module in this pipeline answers the same engineering question at a different
  level of abstraction: \emph{``Given a corrupted measurement, how do we recover the
  truth?''} Distribution math gives us the tools; AWGN and drift models describe the
  enemy; diagnostics identifies which enemy we face; and the Kalman filter, LMS, and
  sensor fusion are three progressively more powerful weapons to fight it.
\end{tcolorbox}

% ====================================================
\section*{References}
% ====================================================
\addcontentsline{toc}{section}{References}

\begin{enumerate}[label={[\arabic*]}]
  \item Kalman, R. E., ``A New Approach to Linear Filtering and Prediction Problems,''
        \textit{Transactions of the ASME -- Journal of Basic Engineering},
        vol.~82, no.~1, pp.~35--45, 1960.

  \item Widrow, B. and Hoff, M. E., ``Adaptive Switching Circuits,''
        \textit{1960 IRE WESCON Convention Record}, pp.~96--104, 1960.

  \item Box, G. E. P. and Muller, M. E., ``A Note on the Generation of Random Normal
        Deviates,'' \textit{The Annals of Mathematical Statistics},
        vol.~29, no.~2, pp.~610--611, 1958.

  \item Papoulis, A. and Pillai, S. U., \textit{Probability, Random Variables, and
        Stochastic Processes}, 4th ed.\ McGraw-Hill, 2002.

  \item Haykin, S., \textit{Adaptive Filter Theory}, 4th ed.\ Prentice Hall, 2002.

  \item Welch, G. and Bishop, G., ``An Introduction to the Kalman Filter,''
        Technical Report TR 95-041, University of North Carolina, 2006.
\end{enumerate}

\end{document}
